% !TEX root = ../main.tex

\documentclass[../main.tex]{subfiles}

% Paper-specific metadata
\title{Placeholder Manuscript Title for Chapter Three}
\manuscriptauthors{Your Name\textsuperscript{a}, Coauthor First\textsuperscript{a}}
\affil{a}{Your Institute, Your Department, Your University, City, Country}

\manuscriptpublished{This manuscript was published by \textit{Journal Name} on Month Day, Year (doi: 10.0000/placeholder).}

\authorcontributions{%
  [Placeholder author contribution statement.]
}

\manuscriptabstract{
  [Placeholder abstract for Chapter 3. This example shows an alternative
  layout where section headings live in \texttt{ch3.tex} and each section's
  body text is a separate subfile -- as opposed to Chapter 2, where each
  subfile defines its own \texttt{\textbackslash section}. Either pattern
  works; pick whichever keeps your files easiest to navigate.]
}

\begin{document}
\manuscriptfront

\section{Introduction}
\input{chapter_3/11_introduction.tex}
\section{Results}
\input{chapter_3/12_results.tex}
\section{Discussion}
\input{chapter_3/13_discussion.tex}
\section{Methods}
\input{chapter_3/14_methods.tex}

\section*{Data, Materials, and Software Availability}
[Placeholder: link to your data/code repository, e.g. Zenodo DOI.]

\section*{Acknowledgements}
[Placeholder acknowledgements for this manuscript.]
\end{document}
